\documentclass[11pt]{article}
\usepackage[utf8]{inputenc}
\usepackage[margin=1in]{geometry}
\usepackage{amsmath}
\usepackage{graphicx}
\usepackage{booktabs}
\usepackage{hyperref}
\usepackage[round]{natbib}
\usepackage{float}

\title{Superficial Superior Colliculus Is Causally Involved in Complex Figure Detection and Encodes Figure--Ground Signals in Mice}
\author{Author A$^{1}$, Author B$^{1}$, Author C$^{1,2}$\\[6pt]
$^{1}$Netherlands Institute for Neuroscience, Royal Netherlands Academy\\
$^{2}$Department of Integrative Neurophysiology, VU University Amsterdam}
\date{}

\begin{document}
\maketitle

\begin{abstract}
Object detection in textured visual scenes requires distinguishing figures from their backgrounds---a computation traditionally attributed to visual cortex. Here we demonstrate that the superficial superior colliculus (sSC) of mice is causally involved in detecting figures defined by complex features on non-homogeneous backgrounds. Bilateral optogenetic inhibition of sSC in GAD2-Cre mice (N=8) significantly decreased performance on all three figure detection tasks---contrast, orientation, and phase---with the strongest effects at early latencies (0--100~ms after stimulus onset). Visually evoked sSC responses were reduced by 33\% during inhibition ($p<0.001$). Electrophysiological recordings in behaving mice (N=5) showed that sSC neurons respond more strongly to figure elements than to identical ground elements for contrast and orientation tasks, and that a linear SVM decoder discriminates figure from ground above chance for all three tasks including phase. Figure--ground discriminability (d-prime) was significantly higher on correct trials than error trials across all tasks, linking sSC neural activity to behavioral decisions. A small group of putative multisensory neurons at deeper locations showed task-related responses peaking at ${\sim}$700~ms, after lick spout presentation but before licking. These findings establish the sSC as a causal contributor to complex figure--ground segregation, expanding its known role beyond simple stimulus detection.
\end{abstract}

\section{Introduction}

Detecting objects in cluttered visual environments is a fundamental perceptual ability. Figure--ground segregation---the neural computation by which local visual elements are classified as belonging to a figure or its background---has been extensively studied in primary visual cortex (V1), where neurons show enhanced responses to figure elements compared to identical stimuli presented in the ground \citep{lamme1995figure, wang2020v1}. However, mice can detect simple stimuli even after V1 silencing or ablation \citep{schneider1969two}, likely through the superior colliculus (SC), a midbrain structure that receives direct retinal input and mediates visually guided behaviors \citep{may2006mammalian, krauzlis2013superior}.

The SC has been implicated in detecting isolated stimuli and simple luminance changes \citep{dean1989visual, boehnke2011colliculus}, and recent work showed that mouse sSC encodes complex visual features independently of V1 input \citep{feinberg2021mouse, hoy2019defined}. However, whether sSC contributes to detecting figures on complex, non-homogeneous textured backgrounds---and whether this contribution is behaviorally relevant---remained unknown.

Here we combine bilateral optogenetic inhibition with electrophysiological recordings during a figure detection task to establish the causal role and neural coding mechanisms of sSC in complex figure--ground segregation.

\section{Results}

\subsection{Optogenetic Inhibition of sSC Impairs Figure Detection}

Eight GAD2-Cre mice \citep{taniguchi2011gad2} received bilateral injections of Cre-dependent ChR2 into sSC, targeting inhibitory neurons. Blue laser activation of these neurons suppressed overall sSC activity. Anesthetized control recordings confirmed that optogenetic inhibition reduced visually evoked sSC responses by 33\% on average ($p<0.001$; Table~\ref{tab:inhibition}), though silencing was not complete.

\begin{table}[H]
\centering
\caption{Optogenetic inhibition effectiveness (anesthetized control, N=3 GAD2-Cre mice).}
\label{tab:inhibition}
\begin{tabular}{lccc}
\toprule
\textbf{Measure} & \textbf{Laser off} & \textbf{Laser on} & \textbf{$p$-value} \\
\midrule
Mean evoked response (sp/s) & $42.3 \pm 5.1$ & $28.3 \pm 4.2$ & $<0.001$ \\
Reduction (\%) & -- & $33.1 \pm 4.8$ & -- \\
Complete silencing & -- & No & -- \\
\bottomrule
\end{tabular}
\end{table}

Mice performed a two-alternative forced choice figure detection task, reporting the position (left vs.\ right) of a figure that differed from the textured background in contrast, orientation, or phase. Inhibition significantly decreased task accuracy for all three figure types (Table~\ref{tab:optobehavior}).

\begin{table}[H]
\centering
\caption{Effect of sSC inhibition on figure detection accuracy (N=8 mice, mean $\pm$ SEM).}
\label{tab:optobehavior}
\begin{tabular}{lcccc}
\toprule
\textbf{Task} & \textbf{No laser (\%)} & \textbf{Laser (\%)} & \textbf{$\Delta$ (\%)} & \textbf{$p$-value} \\
\midrule
Contrast & $82.4 \pm 2.3$ & $71.8 \pm 2.8$ & $-10.6$ & $0.003$ \\
Orientation & $78.1 \pm 2.7$ & $68.3 \pm 3.1$ & $-9.8$ & $0.006$ \\
Phase & $74.6 \pm 3.0$ & $65.2 \pm 3.4$ & $-9.4$ & $0.011$ \\
\bottomrule
\end{tabular}
\end{table}

\subsection{Temporal Profile of Inhibition Effects}

Inhibition was most effective when applied at early latencies (0--100~ms after stimulus onset), with diminished effects at later time points (Table~\ref{tab:temporal}).

\begin{table}[H]
\centering
\caption{Temporal profile of sSC inhibition effects on accuracy ($\Delta$\% from no-laser baseline).}
\label{tab:temporal}
\begin{tabular}{lccc}
\toprule
\textbf{Latency window} & \textbf{Contrast} & \textbf{Orientation} & \textbf{Phase} \\
\midrule
0--100 ms (early) & $-12.4 \pm 2.1$ & $-11.8 \pm 2.4$ & $-10.2 \pm 2.6$ \\
100--200 ms (late) & $-6.3 \pm 1.8$ & $-5.9 \pm 2.0$ & $-5.1 \pm 2.2$ \\
\bottomrule
\end{tabular}
\end{table}

\subsection{sSC Neurons Encode Figure--Ground Differences}

Electrophysiological recordings (N=5 C57BL/6J mice) using silicon probes (Neuropixels and 32-channel linear arrays; \citealt{jun2017neuropixels}) revealed that neurons with receptive fields (RFs) inside the figure responded more strongly when their RF contained figure elements versus identical ground elements (Table~\ref{tab:figground}).

\begin{table}[H]
\centering
\caption{Population figure--ground response differences (RF inside figure, cluster-based permutation test).}
\label{tab:figground}
\begin{tabular}{lccc}
\toprule
\textbf{Task} & \textbf{Figure $>$ Ground?} & \textbf{Significant clusters} & \textbf{$p$-value} \\
\midrule
Contrast & Yes & Multiple time windows & $<0.05$ \\
Orientation & Yes & Multiple time windows & $<0.05$ \\
Phase & Not at population level & None & $>0.05$ \\
\bottomrule
\end{tabular}
\end{table}

Neurons with RFs on the figure edge also showed figure--ground modulation, with significant time clusters for contrast and orientation tasks.

\subsection{SVM Decoding of Figure Versus Ground}

A linear SVM classifier with a 50~ms sliding window (10~ms steps) decoded figure versus ground identity from sSC population activity above chance for all three tasks, including phase (Table~\ref{tab:svm}).

\begin{table}[H]
\centering
\caption{SVM decoding performance (peak accuracy, 50~ms window).}
\label{tab:svm}
\begin{tabular}{lcccc}
\toprule
\textbf{Task} & \textbf{Peak accuracy (\%)} & \textbf{Chance (\%)} & \textbf{Significant bins} & \textbf{$p$} \\
\midrule
Contrast & $72.4$ & $50$ & $14$ & $<0.01$ \\
Orientation & $68.7$ & $50$ & $11$ & $<0.01$ \\
Phase & $61.3$ & $50$ & $7$ & $<0.05$ \\
\bottomrule
\end{tabular}
\end{table}

The successful decoding of phase-defined figures---despite the absence of a significant population-level response difference---indicates that consistent but subtle single-neuron response differences carry sufficient information for figure--ground discrimination.

\subsection{Neural Discriminability Predicts Behavioral Outcome}

Figure--ground discriminability (d-prime from ROC analysis; \citealt{green1966signal}) was significantly higher on correct trials (hits) than error trials across all tasks (Table~\ref{tab:dprime}), demonstrating that sSC activity is behaviorally relevant.

\begin{table}[H]
\centering
\caption{Figure--ground discriminability (d-prime) on hit versus error trials.}
\label{tab:dprime}
\begin{tabular}{lcccc}
\toprule
\textbf{Task} & \textbf{d$'$ hits} & \textbf{d$'$ errors} & \textbf{$\Delta$d$'$} & \textbf{$p$-value} \\
\midrule
Contrast & $0.84 \pm 0.12$ & $0.31 \pm 0.09$ & $0.53$ & $0.002$ \\
Orientation & $0.72 \pm 0.11$ & $0.28 \pm 0.08$ & $0.44$ & $0.005$ \\
Phase & $0.48 \pm 0.10$ & $0.18 \pm 0.07$ & $0.30$ & $0.018$ \\
\bottomrule
\end{tabular}
\end{table}

\subsection{Putative Multisensory Neurons Show Late Task-Related Responses}

A small group of neurons recorded at slightly deeper locations showed peak responses at ${\sim}$700~ms---after lick spout presentation but before licking---with figure--ground modulation similar to visual neurons but at a much later latency (Table~\ref{tab:multisensory}).

\begin{table}[H]
\centering
\caption{Response properties of visual versus putative multisensory SC neurons.}
\label{tab:multisensory}
\begin{tabular}{lcc}
\toprule
\textbf{Property} & \textbf{Visual neurons} & \textbf{Putative multisensory} \\
\midrule
Peak response time (ms) & $80$--$150$ & ${\sim}700$ \\
Figure $>$ Ground & Yes & Yes (trend) \\
Depth & sSC (superficial) & Slightly deeper \\
Task modulation & Visual stimulus-locked & Decision/motor-locked \\
\bottomrule
\end{tabular}
\end{table}

\section{Discussion}

This study establishes three key findings about the role of the superficial superior colliculus in complex visual processing. First, bilateral sSC inhibition causally impairs detection of figures defined by contrast, orientation, and phase on textured backgrounds. This extends the known role of SC from simple stimulus detection \citep{schneider1969two, dean1989visual} to complex figure--ground segregation---a computation previously attributed primarily to V1 \citep{lamme1995figure, wang2020v1}.

Second, the early temporal window (0--100~ms) during which inhibition is most effective is consistent with the sSC contributing to feedforward visual processing \citep{resulaj2018first}. The diminished effect at later latencies (100--200~ms) may reflect compensatory processing by cortical circuits.

Third, the correlation between sSC neural discriminability and behavioral outcome provides direct evidence that sSC figure--ground signals contribute to perceptual decisions. The successful decoding of phase-defined figures---despite weak population-level effects---suggests that sparse single-neuron signals in sSC carry behaviorally relevant information.

The putative multisensory neurons with late responses may reflect SC's known role in integrating visual and motor/decision signals \citep{wallace2004multisensory, krauzlis2013superior}, suggesting that SC contributes not only to sensory processing but also to the decision stage of figure detection.

\section{Methods}

\subsection{Animals}
C57BL/6J mice (N=5, male, 2--5 months) were used for electrophysiology. GAD2-Cre mice (N=8, 6 male / 2 female; Jackson \#028867) were used for optogenetics. All were housed on a reversed 12h light/dark cycle.

\subsection{Optogenetic Inhibition}
Cre-dependent ChR2 (ssAAV-9/2-hEF1a-dlox-hChR2(H134R)\_mCherry, titer $5.4 \times 10^{12}$) was injected bilaterally into sSC. Optic fibers were implanted. Blue laser activated inhibitory neurons. A blue LED flashed randomly above the head to prevent behavioral cueing.

\subsection{Behavioral Task}
Mice performed a 2-AFC figure detection task (left/right lick on Y-shaped spout). Figures differed from textured backgrounds in contrast, orientation, or phase. Response was allowed from 200~ms post-stimulus. Mice ran 100--250 trials/session over 2--5 months.

\subsection{Electrophysiology}
Silicon probes (Neuropixels, 32-channel linear arrays) targeted sSC. Neurons were classified by RF location (inside figure, edge, outside). Histological verification used DiI + DAPI staining.

\subsection{Analysis}
Population responses used cluster-based permutation tests \citep{maris2007nonparametric}. SVM classification used 50~ms sliding windows (10~ms steps). ROC-derived d-prime compared hit and error trials \citep{green1966signal}.

\section{Conclusion}

The superficial superior colliculus is causally involved in detecting figures defined by complex features on textured backgrounds, encodes figure--ground signals that predict behavioral outcomes, and contains a population of putative multisensory neurons with decision-time responses. These findings establish the sSC as an integral component of the neural circuit for complex visual scene analysis.

\bibliographystyle{plainnat}
\bibliography{references}

\end{document}
